\section{Name Matching Details}
\label{app:matching_details}

The usage analysis in RQ1 requires joining two independently-maintained name
spaces: OpenRouter's routing slugs (347 distinct \texttt{permaslug}s appear in
the daily top-50 between Jan.\ 2025 and Jun.\ 2026) and LMArena's model names
(364 distinct names in the same window). The two disagree on provider prefixes
(\texttt{anthropic/}), on routing-variant suffixes (\texttt{:free},
\texttt{:nitro}), on whether a release carries a date stamp and in what format
(\texttt{-20260217} vs.\ \texttt{-2026-02-17} vs.\ nothing), and on
preview-to-GA renames. We match in three layers, each of which either produces a
unique match or abstains.

\paragraph{Layer 1: canonical token multiset.} We lowercase the slug, drop
everything before the last \texttt{/} and everything after a \texttt{:}, and
split the remainder into maximal alphabetic and numeric runs. Two names match if
their sorted multisets of alphabetic tokens \emph{and} of numeric tokens agree,
which makes the comparison insensitive to component order
(\texttt{qwen3-32b-instruct} vs.\ \texttt{qwen3-instruct-32b}) and to separator
style. A match is accepted only when exactly one Arena name shares the key;
ambiguous keys are left unmatched. This yields 89 matches.

\paragraph{Layer 2: date-stamp deletion.} We repeat Layer 1 after deleting
numeric tokens that are exactly 8 digits long, i.e.\ genuine
\texttt{YYYYMMDD} stamps, again requiring uniqueness. This yields a further 30
matches. We deliberately stop here rather than applying general fuzzy or
trailing-token trimming: more permissive rules silently produce wrong matches
(trimming trailing numerics collapses \texttt{kimi-k2.5} into \texttt{kimi-k2}),
and a wrong match is far more damaging here than an abstention, since it
attributes real token volume to the wrong capability profile.

\paragraph{Layer 3: hand-verified overrides.} Every remaining OpenRouter model
with material token volume was read and either linked by hand or explicitly left
unmatched. The override table has 60 entries, keyed on the base slug so that
routing variants of the same model inherit the mapping; it resolves 73 of the
347 slugs. Entries are tagged \emph{high} confidence (46 entries; pure formatting
or token-order differences, or a single obvious date-dropped candidate, e.g.\
\texttt{grok-4-07-09}$\to$\texttt{grok-4-0709}) or \emph{medium} (14 entries;
a genuine judgment call, such as mapping a bare undated alias to the more recent
of two dated Arena checkpoints, choosing among reasoning-effort variants of the
same release, or picking the reference precision variant when OpenRouter does
not distinguish bf16 from fp8).

\paragraph{Deliberate exclusions.} Some OpenRouter models have no defensible
Arena counterpart and are left unmatched by design rather than by failure:
OpenRouter's stealth/cloaked pre-release codenames (\texttt{owl-alpha},
\texttt{horizon-beta}, \texttt{quasar-alpha}, and roughly fifteen others), for
which no reliable public mapping to the eventual release name exists; embedding
models; vision-only, safety, base, and distilled checkpoints that are not
separately rated in the Arena text categories; and niche community fine-tunes
never submitted to the Arena.

\paragraph{Coverage and treatment of the remainder.} The three layers together
match 192 of 347 slugs (89 exact, 30 after date deletion, 73 override; of the
override matches, 56 are high- and 17 medium-confidence), covering 83.5\% of
pooled token volume and a mean of 79.2\% of monthly volume, ranging from 60.7\%
to 91.5\% across the 18 months. Coverage dips when a heavily-used release is new
enough not to be aliased yet; the largest single unmatched models are
\texttt{grok-code-fast-1} (4.1\% of pooled tokens),
\texttt{gemini-2.5-flash-lite} (2.4\%), and the stealth slug
\texttt{owl-alpha} (1.0\%). Throughout, unmatched volume is treated as
\emph{not} lying within any threshold of the top-rated model, so the
near-frontier usage shares we report are a conservative floor on the true
concentration of usage near the frontier, and the reported share curve cannot
exceed the matched-coverage ceiling by construction.

\section{Open Source Models}
\label{app:open_source_models}

The scaling analysis in RQ2 needs capability scores and an observable measure of
investment for the same models. We take scores from the Open LLM Leaderboard's
per-model results (3,861 leaderboard rows) on its six objective capability
dimensions --- IFEval, BBH, MATH Level 5, GPQA, MuSR, and MMLU-Pro --- and
parameter counts from the associated Hugging Face repository metadata. We drop
rows with missing or zero parameter counts (10 rows) and rows the leaderboard
flags as contaminated or as weight merges (none survive into our sample in this
snapshot). Each of the six dimensions is standardized to zero mean and unit
variance across this full 3,851-model pool \emph{before} any lab restriction, so
that the scale of a slope is set by the leaderboard as a whole rather than by
the subsample we ultimately fit on; without this, the six panels of
Figure~\ref{fig:scaling_curves_normalized} would not be comparable, since the
underlying benchmarks have very different raw ranges.

\paragraph{Which organizations.} The leaderboard contains 650 distinct uploading
organizations, the large majority of which are community fine-tunes of someone
else's base model. For those, parameter count reflects the technology of the
base-model provider rather than of the uploader, so a fitted scaling slope would
not be that organization's production function. We therefore restrict to 17
hand-identified organizations that pre-train their own models, and then require
a genuine scaling ladder to fit against: at least 7 models, spanning at least
$10^{0.6}\approx4\times$ in parameter count, at at least 4 distinct sizes.
Table~\ref{tab:ollm_lab_screen} reports the screen. Fifteen labs survive,
covering 293 leaderboard rows over 278 distinct repositories (a repository
re-submitted under a later revision contributes more than one row; we keep both
as independent evaluation runs, and all cross-validation splits are by
repository so that duplicate rows never straddle a train/test boundary). The
retained models span 0.08B to 140.6B parameters and upload dates from March 2022
to March 2025, and are 136 instruction-tuned, 113 base, and 44 other
checkpoints; all slope fits include an instruction-tuned indicator so that a
lab's mix of base and post-trained releases cannot masquerade as a scaling
response.

\begin{table}[t]
    \centering
    \begin{tabular}{lccccc}
        \toprule
        Lab & Models & Rows & Sizes & Param.\ range (B) & Span (dex) \\
        \midrule
        Qwen           & 58 & 60 & 22 & 0.49--111.2 & 2.35 \\
        Google         & 29 & 32 & 17 & 0.08--27.2  & 2.55 \\
        IBM Granite    & 20 & 20 &  7 & 1.33--8.2   & 0.79 \\
        Meta           & 19 & 20 &  7 & 1.24--70.6  & 1.76 \\
        NousResearch   & 19 & 19 & 12 & 3.21--70.6  & 1.34 \\
        Microsoft      & 16 & 18 & 11 & 0.34--42.0  & 2.09 \\
        Mistral        & 16 & 18 & 11 & 7.24--140.6 & 1.29 \\
        NVIDIA         & 18 & 18 & 13 & 1.52--72.7  & 1.68 \\
        AI2            & 14 & 16 &  9 & 1.18--70.6  & 1.78 \\
        TII            & 16 & 16 &  8 & 1.67--40.0  & 1.38 \\
        HuggingFaceTB  & 12 & 14 &  6 & 0.13--1.7   & 1.12 \\
        EleutherAI     & 12 & 12 & 12 & 0.15--20.7  & 2.14 \\
        Arcee AI       & 10 & 11 &  8 & 3.09--72.7  & 1.37 \\
        DeepSeek       & 11 & 11 &  9 & 1.78--70.6  & 1.60 \\
        Cohere For AI  &  8 &  8 &  4 & 8.03--103.8 & 1.11 \\
        \midrule
        \textit{01.AI}       & \textit{19} & \textit{19} & \textit{3} & \textit{6.06--34.4} & \textit{0.75} \\
        \textit{BigScience}  &  \textit{5} &  \textit{5} & \textit{5} & \textit{0.56--7.1}  & \textit{1.10} \\
        \bottomrule
    \end{tabular}
    \caption{The open-weight lab screen. ``Models'' counts distinct Hugging Face
    repositories and ``Rows'' counts leaderboard submissions; ``Sizes'' is the
    number of distinct parameter counts and ``Span'' is
    $\log_{10}(\max/\min)$ of parameter count. The two italicized labs fail the
    screen: 01.AI has ample models but only 3 distinct sizes over a $0.75$-dex
    range, and BigScience has too few models, so neither supports a per-lab
    scaling fit.}
    \label{tab:ollm_lab_screen}
\end{table}

\paragraph{Limitations.} Three caveats bound what this sample can identify.
First, parameter count is only a partial proxy for investment: it omits the
training token budget, so two labs at the same parameter count may sit at very
different points on the compute frontier, and a lab whose releases are
systematically over- or under-trained relative to the pool will have its slope
mis-attributed to technology rather than to token budget. Second, the sample is
open-weight by necessity --- parameter counts are simply not published for the
closed frontier flagships that drive RQ1 --- so the labs here are not the same
labs, and the estimated scaling behavior should not be read as a direct estimate
of the frontier firms' production functions. Third, two of the fifteen retained
labs (Arcee AI and NousResearch) primarily fine-tune others' base models rather
than pre-train, so their fitted scaling response is partly derivative of the
base-model providers'; we retain them because they pass the ladder screen and
excluding them does not change the qualitative conclusions, but their per-lab
slopes should be read with that in mind.
